
\documentclass[11pt]{amsart}

\usepackage{amsmath}
\usepackage{amsfonts}
\usepackage{amssymb}
\usepackage{amsthm}
\usepackage{graphicx}


\title{A non subdirectly irreducible algebra generating a variety of DMM that has JEP}

\begin{document}

\maketitle

\allowdisplaybreaks

%The following is a non-simple algebra generating a variety with JEP. It is essentially the reflection of $\mathbf{2}$ with an idempotent added in the middle.
%
%This algebra has three subalgebras, $C_4 = \{0,2,4,6\}$, $\text{R}(2) = \{0,1,2,4,5,6\}$ and $\{0,2,3,4,6\}$. It has one interesting congruence which identifies $1-2$, and $4-5$.
%
%\input{output/E-7-1}
%\pagebreak[4]

The purpose of this document is to provide a counter example to the following conjecture.

\begin{quotation}
\textbf{Conjecture.}\, Let $\mathsf{K}$ be a variety of DMMs with the JEP, that is not PSC.  Then $\mathsf{K} = \mathbb{Q}(\mathbf{A})$ for some subdirectly irreducible algebra $\mathbf{A}$.
\end{quotation}

Were this true, it would provide a refinement of the following known characterisation.

\begin{quotation}
\textbf{Theorem.}\, A congruence distributive variety $\mathsf{K}$ has the JEP iff $\mathsf{K} = \mathbb{V}(\mathbf{A})$ for some algebra $\mathbf{A}$ such that $(\mathbb{HSP}_u(\mathbf{A}))_{si}$ is contained in $\mathbb{ISP}_u(\mathbf{A})$.
\end{quotation}

Notice that when $\mathbf{A}$ is finite we can remove the $\mathbb{P}_u$ operators above.

To construct a counter-example to the conjecture above we will exhibit an algebra which satisfies the condition of the theorem, but can not be obtained from a single subdirectly irreducible.

%We know from previous work that this algebra must have $\mathbf{C}_4$ as a subalgebra. It should not have $\mathbf{C}_4$ as a homomorphic image, however, otherwise it might generate a PSC variety. To make sure $\mathbf{C}_4$ is not a homomorphic image, we will arrange that our algebra has an odd number of elements, so that there is an element which is equal to its own negation. We will also insist that this element is idempotent. So that the following algebra, I'll call $\mathbf{E}$, is a subalgebra.
Consider the following algebra I'll call $\mathbf{E}$.
{\tiny
\thicklines
\begin{center}
\begin{picture}(50,88)(120,20)
%
%%%%%%
%
\put(122,45.5){\circle*{4}}
\put(122,63){\circle*{4}}
\put(122,81.5){\circle*{4}}
\put(122,98){\circle*{4}}
\put(122,78){\line(0,1){20}}
\put(122,28){\circle*{4}}
\put(122,28){\line(0,1){60}}

\put(127,96){\small $f^2=f \cdot a$}

\put(127,79){\small $f$}
\put(127,61){\small{$a=a^2=\neg a$}}
\put(127,43){\small $e$}

\put(127,25){\small $\neg(f^2)$}
\end{picture}
\end{center}
}

Let $\mathbf{2}$ be the 2-element boolean algebra and $\mathbf{S_3}$ the 3-element Sugihara monoid. Recall that one can take reflections of $\mathbf{2}$ and $\mathbf{S_3}$ to obtain De Morgan monoids $\text{R}(\mathbf{2})$ and $\text{R}(\mathbf{S_3})$. Recall that $\mathbf{C_4}$ is the only interesting subalgebra and homomorphic image of these reflections, therefore, up to isomorphism,
$$ (\mathbb{HSP}_u(\text{R}(\mathbf{2})))_{si} = \{\text{R}(\mathbf{2}),\mathbf{C_4}\} \quad (\mathbb{HSP}_u(\text{R}(\mathbf{S_3})))_{si} = \{\text{R}(\mathbf{S_3}),\mathbf{C_4}\}. $$
One can insert a new idempotent element $a$ in the middle of these reflections, so that $a = \neg a$, and $a \cdot x = f^2$ for all $x > a$, and $a \cdot x = a$ for all $\bot < x \leq a$. When we do this we obtain De Morgan monoids called $\mathbf{E-7-1}$ and $\mathbf{E-9-1}$ (for no reason other than that is what their uacalc files are called on my computer). They each have $\mathbf{E}$ their only interesting homomorphic image. Up to isomorphism,
$$ (\mathbb{HSP}_u(\mathbf{E-7-1}))_{si} = \{\mathbf{E-7-1},\text{R}(\mathbf{2}),\mathbf{E},\mathbf{C_4}\}$$
\begin{equation}
\label{eq:9si}
(\mathbb{HSP}_u(\mathbf{E-9-1}))_{si} = \{\mathbf{E-9-1},\text{R}(\mathbf{S_3}),\mathbf{E},\mathbf{C_4}\}.
\end{equation}



To construct our algebra, we use a similar tactic. It is constructed by adding an idempotent element (7) to the reflection of $\mathbf{2} \times \mathbf{S_3}$.

\input{output/E-15-1}
%\pagebreak[4]

$\mathbf{E-15-1}$ has the following subalgebras.
\begin{align*}
C_4 &= \{0,4,10,14\} \\
\text{R}(2) &= \{0,3,4,10,11,14\} \\
\text{R}(S_3) &= \{0,2,4,6,8,10,12,14\} \\
E &= \{0,4,7,10,14\} \\
E-7-1 &= \{0,3,4,7,10,11,14\} \\
E-9-1 &= \{0,2,4,6,7,8,10,12,14\} \\
\text{R}(2 \times S_3) &= \{0,1,2,3,4,6,8,9,10,11,12,13,14\}
\end{align*}

It has the following interesting congruences.
\begin{align*}
E-9-1:&\{0\},\{1,2\},\{3,4\},\{5,6\},\{7\},\{8,9\},\{10,11\},\{12,13\},\{14\} \\
E-7-1:&\{0\},\{1,3,5\},\{2,4,6\},\{7\},\{8,10,12\},\{9,11,13\},\{14\} \\
E:&\{0\},\{1,2,3,4,5,6\},\{7\},\{8,9,10,11,12,13\},\{14\}
\end{align*}
Note that $(\mathbb{HSP}_u(\mathbf{E-15-1}))_{si} \subseteq \mathbb{ISP}_u(\mathbf{E-15-1})$. Therefore $\mathbb{V}(\mathbf{E-15-1})$ has the JEP by the theorem above.

Let us now prove that $\mathsf{K} = \mathbb{V}(\mathbf{E-15-1}) \neq \mathbf{V}(\mathbf{A})$ for any subdirectly irreducible $\mathbf{A}$. Suppose not, then, by Jonsson's theorem,
$$\mathbf{A} \in (\mathbb{HSP}_u(\mathbf{E-15-1}))_{si} = \{\mathbf{E-9-1}, \mathbf{E-7-1}, \text{R}(\mathbf{S_3}),\text{R}(\mathbf{2}),\mathbf{E},\mathbf{C_4}\}.$$
If $\mathbf{A}$ is any of the algebras above apart from $\mathbf{E-9-1}$, then
$$\mathbf{E-9-1} \in \mathsf{K}_{si} \subseteq \mathbb{HSP}_u(\mathbf{A}),$$
which is false on cardinality grounds. Therefore, $\mathbf{A} = \mathbf{E-9-1}$. But notice, from (\ref{eq:9si}), that $\mathbf{E-7-1} \notin (\mathbb{HSP}_u(\mathbf{E-9-1}))_{si}$, which contradicts the assumption that $\mathsf{K} = \mathbb{V}(\mathbf{A})$.

Lastly notice that since $\mathsf{K}$ contains $\mathbf{E}$, which does not map onto $\mathbf{C_4}$, we know that $\mathsf{K} \notin \mathsf{M}$, in particular $\mathsf{K}$ is not PSC.

We have therefore provided an example of a variety of DMMs with the JEP that is not PSC, but which is not generated as a variety (nevermind a quasivariety) by a single subdirectly irreducible algebra.
\end{document}
